\let\negmedspace\undefined
\let\negthickspace\undefined
\documentclass[journal]{IEEEtran}
\usepackage[a4paper, margin=10mm, onecolumn]{geometry}
%\usepackage{lmodern} % Ensure lmodern is loaded for pdflatex
\usepackage{tfrupee} % Include tfrupee package
\setlength{\headheight}{1cm} % Set the height of the header box
\setlength{\headsep}{0mm}  % Set the distance between the header box and the top of the text
\usepackage{gvv-book}
\usepackage{gvv}
\usepackage{cite}
\usepackage{amsmath,amssymb,amsfonts,amsthm}
\usepackage{algorithmic}
\usepackage{graphicx}
\usepackage{float}
\usepackage{textcomp}
\usepackage{xcolor}
\usepackage{txfonts}
\usepackage{listings}
\usepackage{enumitem}
\usepackage{mathtools}
\usepackage{gensymb}
\usepackage{comment}
\usepackage[breaklinks=true]{hyperref}
\usepackage{tkz-euclide} 
\usepackage{listings}
% \usepackage{gvv}                                        
\def\inputGnumericTable{}                                 
\usepackage[latin1]{inputenc}                                
\usepackage{color}                                            
\usepackage{array}                                            
\usepackage{longtable}                                       
\usepackage{calc}                                             
\usepackage{multirow}                                         
\usepackage{hhline}                                           
\usepackage{ifthen}                                           
\usepackage{lscape}
\usepackage{tikz}
\usetikzlibrary{patterns}
\begin{document}
\bibliographystyle{IEEEtran}
\vspace{3cm}
\title{4.11.31}
\author{ee25btech11054-S.Harsha Vardhan Reddy}
\maketitle
% \maketitle
% \newpage
% \bigskip
{\let\newpage\relax\maketitle}
\renewcommand{\thefigure}{\theenumi}
\renewcommand{\thetable}{\theenumi}
\setlength{\intextsep}{10pt} % Space between text and floats
\numberwithin{equation}{enumi}
\numberwithin{figure}{enumi}
\renewcommand{\thetable}{\theenumi}

\textbf{Question}\\
Find the area of the region bounded by line $y=3x+2$, the X axis and the ordinates $x=-2$ and $x=1$.\\

\textbf{Solution:}\\
Let
\begin{align}
    \vec{A}=\myvec{-2\\0}\\
    \vec{C}=\myvec{1\\0}
\end{align}
Let $\vec{D}$ and $\vec{E}$ be the vectors on the line corresponding to $x= -2$ and $x=1$.\\

Given line equation is 
\begin{align}
    -3x+y=2
\end{align}
which can be expressed as
\begin{align}
    \vec{n}^T\vec{x}=c\\
    \implies \myvec{-3 & 1}\myvec{x\\y}=2
\end{align}
\begin{align}
    \vec{n}=\myvec{-3\\1} \text{ and } \vec{x}=\myvec{x\\y} \text{ and } c=2
\end{align}
Let us find the vector $\vec{D}$:
\begin{align}
    \myvec{-3 & 1}\myvec{-2\\y}=2 \implies y=-4\\
    \implies \vec{D}= \myvec{-2\\-4}
\end{align}
As $y<0$, we should find the $\vec{B}$ where the line meets the $x$ axis:
\begin{align}
    \myvec{-3 & 1}\myvec{x\\0}=2 \implies 3x=-2\\
    \implies \vec{B}= \myvec{-\frac{2}{3}\\0}
\end{align}
Let us find the vector $\vec{E}$:
\begin{align}
    \myvec{-3 & 1}\myvec{1\\y}=2 \implies y=5\\
    \implies \vec{E}= \myvec{1\\5}
\end{align}
The area to be computed is area of $\triangle EBC$ + area of $\triangle ABD$:
\begin{align}
    \text{ar}(\triangle ABD)=\frac{1}{2}\norm{\brak{\vec{A}-\vec{B}}\times\brak{\vec{A}-\vec{D}}}\\
    =\frac{1}{2}\norm{\myvec{-\frac{4}{3}\\0}\times\myvec{0\\4}}=\frac{8}{3}
\end{align}
\begin{align}
    \text{ar}(\triangle EBC)=\frac{1}{2}\norm{\brak{\vec{E}-\vec{C}}\times\brak{\vec{B}-\vec{C}}}\\
    =\frac{1}{2}\norm{\myvec{0\\5}\times\myvec{-\frac{5}{3}\\0}}=\frac{25}{6}
\end{align}
\begin{align}
    \implies \text{Area of the region} =\frac{8}{3}+\frac{25}{6}=\frac{41}{6}
\end{align}

\begin{figure}[H]
    \centering
    \includegraphics[width=0.96\columnwidth]{figs/01.png}
    \caption{Region bounded by the line and ordinates}
    \label{fig-1}
\end{figure}
\end{document}